\section{Introduction}

Wireless routing is often taught in the context of connected mobile ad hoc networks, where a multi-hop end-to-end path is assumed to exist long enough for route discovery and forwarding. Intermittently connected mobile networks (ICMNs) violate that assumption. In sparse or highly mobile settings, nodes meet only occasionally, and a source may never have a contemporaneous path to the destination. Delay-tolerant networking addresses this setting through store-carry-forward communication: a node buffers data, carries it while moving, and forwards it opportunistically when it encounters another node~\cite{fall2003dtn}.

This project studies a simple but important question: how much controlled replication is needed to make opportunistic delivery effective in a sparse wireless network? Uncontrolled flooding approaches such as Epidemic Routing can achieve strong delivery performance, but at the cost of excessive transmissions and resource use~\cite{vahdat2000epidemic}. More broadly, DTN routing has explored several competing approaches, including oracle- or knowledge-based forwarding~\cite{jain2004routing}, history-based probabilistic forwarding such as PRoPHET~\cite{lindgren2012prophet}, and prioritization-oriented schemes such as MaxProp~\cite{burgess2006maxprop}. Binary Spray and Wait was proposed as a middle ground: rather than flood the network, a source injects a bounded number of copies and lets those copies spread in a disciplined way~\cite{spyropoulos2005spraywait,spyropoulos2008multiplecopy}. This makes it a strong fit for a small teaching simulator because it exposes a real design tradeoff between reliability, delay, and overhead without requiring heavy protocol machinery.

We implement Binary Spray and Wait in PyWiSim, a minimal event-driven wireless simulator intended for rapid protocol prototyping. PyWiSim is especially suitable for this project because it already includes an encounter-based extension for DTN-style pairwise contacts, allowing the protocol logic to remain compact while still operating over an explicitly wireless abstraction. Our goal is not to claim a production-grade DTN evaluation, but to produce a careful miniature research study: we design the protocol, instrument it to separate data and control transmissions, and evaluate how performance changes with network size and copy budget.

The main contributions of this report are:
\begin{itemize}
    \item An implementation of Binary Spray and Wait over PyWiSim's encounter-based ICMN model.
    \item A reproducible evaluation script that varies both network size and copy budget over 50 independent trials per setting.
    \item An analysis showing that larger copy budgets substantially improve delivery probability in larger networks, while increasing data-forwarding cost and reducing control overhead by completing delivery sooner.
\end{itemize}
